% Compile from this directory with: pdflatex project_10.tex (twice).
\documentclass[11pt,a4paper]{article}
\usepackage[T1]{fontenc}
\usepackage[utf8]{inputenc}
\usepackage[margin=22mm]{geometry}
\usepackage{lmodern}
\usepackage{amsmath,amssymb}
\usepackage{enumitem}
\usepackage{xcolor}
\usepackage{microtype}
\usepackage{hyperref}
\definecolor{epflred}{HTML}{E3000B}
\hypersetup{colorlinks=true,urlcolor=epflred,linkcolor=epflred,
  pdftitle={ENG-654 Project 10: Branch-Dependent Planning with the PUMA Positioning Arm}}
\setlength{\parindent}{0pt}
\setlength{\parskip}{6pt}
\setlist[itemize]{leftmargin=*,itemsep=5pt,topsep=4pt}
\urlstyle{tt}

\begin{document}
{\small\textbf{EPFL \textbar\ ENG-654}\hfill Kinematics-Grounded Motion Planning for Robots}
\vspace{5mm}

{\color{epflred}\Large\bfseries Project 10}

{\LARGE\bfseries Branch-Dependent Planning with the PUMA Positioning Arm\par}

\textbf{Format:} Group of two students; simulation-based project.\\
\textbf{Course foundations:} Lectures 3--7; Exercise 3.

\section*{Motivation}
A wrist-center position may be reached by several shoulder/elbow postures.
Those alternatives can have different continuations along the same Cartesian
curve. Restricting PUMA to its three-joint positioning task makes the global
singularity structure visible while retaining the planning question: which
initial solution can follow the complete path?

\section*{Task}
\textbf{Overall objectives.} Treat PUMA's first three joints as a noncuspidal 3R wrist-center positioning
robot and explain path feasibility through reduced singularity maps and the
distribution of its IK solutions.

\begin{itemize}
  \item \textbf{Extract the position task.} Identify the wrist center from the PUMA
axis geometry and derive its position as a function of $q_1,q_2,q_3$.
Validate the D--H and URDF-based descriptions. The wrist joints are irrelevant
to this point's position; do not replace the wrist center with the tool tip.

  \item \textbf{Build the global reference maps.} Derive $J_p$ and plot $\det J_p$
and its zero contours in $(q_2,q_3)$. Map the singular curves into $(\rho,z)$
and count distinct arm IKs. Identify two-/four-IK regions where present and
explain if this specific geometry lacks one of those regions.

  \item \textbf{Select and classify initial solutions.} Choose one regular loop and
two translated or scaled variants approaching different arm singularities.
Enumerate the initial arm IKs and identify their shoulder/elbow categories.
Overlay them on the reduced plot and record mathematical versus
joint-limit-admissible configurations.

  \item \textbf{Continue every admissible branch.} Follow each workspace curve from
each valid initial arm configuration. Use warm-started continuation and
analytical enumeration as a cross-check. Refine the step size near critical
events; preserve branch identity instead of replacing a failed solution with
an unrelated IK at the next sample.

  \item \textbf{Explain completion and failure.} Compare completion fraction,
joint travel, joint-limit margin, and the smallest singular value of $J_p$.
Relate each failed lift to its configuration-space separator and workspace
critical curve. A workspace point can have another regular IK even when the
selected continuation fails.
\end{itemize}

\newpage
\section*{Specifics and scope}
Use one PUMA positioning model and three position curves, with about
200 samples per curve before local refinement. Keep orientation outside the
task and use the same joint limits and tolerances for every comparison.
Explore mathematical branch structure first, then report limit filtering.
The goal is an explanation of branch-dependent feasibility, not a new general
path-search algorithm.

Use the position Jacobian $J_p$ for the 3R task.
Confirm that the meridian reduction is valid for the selected model and
operational point. Count distinct real IKs modulo periodicity, then filter
joint limits separately. Identify two-/four-IK regions where present and
explain any absent count rather than forcing both types into every map. Show both determinant values and the zero contour,
and distinguish sampled connectivity from a proof. Refine the decisive
region instead of claiming completeness from a coarse grid.

\section*{Expected outcome}
Understand the PUMA arm's D--H model, noncuspidal branch separation, and
the difference between workspace reachability and a feasible lift. Deliver
reduced determinant maps with overlaid IKs and trajectories, workspace IK-count
maps, and a branch-by-branch path comparison.

Submit reproducible code, model parameters and test data, the required plots,
a concise 5--6-page report, and a short simulation demonstration. Explain
both successful cases and failures; conclusions must follow the numerical
and geometric evidence rather than an expected outcome.

\section*{Resources}
The PUMA 560 URDF is
\path{lectures_main/assets/models/puma/puma560_robot.urdf}; its STL meshes are
in the same directory. The lecture website provides axis/frame inspection,
IK visualization, and path demonstrations. Confirm joint offsets and the
selected terminal frame when importing the model.
The Lecture 7 PUMA demonstration already follows wrist-center position with
fixed wrist joints and is directly relevant to this project. Reuse its
continuation and visualization as a starting point, checking the analytical
arm IK against Lectures 3--4.

Use the lecture website to verify D--H parameters, visualize IKs and
singularities, and simulate paths. All listed paths are relative to the
repository root. Start the website following \path{lectures_main/README.md}.

\section*{Workload and collaboration}
Budget approximately 60 hours per student (120 person-hours per pair)
for this 2-ECTS project, following EPFL's 30-hours-per-credit convention.\footnote{\href{https://www.epfl.ch/education/teaching/teaching-guide-2/getting-started/design-a-course_1/}{EPFL: Design a Course, workload guidance}.}
Deduct any lectures or exercises included in those credits. Allocate roughly
20 team-hours to model validation, 50 to the core work, 30 to experiments,
and 20 to reporting. Reuse the specified IK and simulations. Share validation
and interpretation even if modelling and implementation are divided between
students. Hardware execution is outside scope.

\end{document}
